\unirsection{A fondo}

Los siguientes libros pueden servir de material de apoyo y para profundizar más en los contenidos de este tema.

\textbf{Nielsen, M. A., Chuang, I. L. \textit{Quantum Computation and Quantum Information}. Cambridge University Press, 2010}

El capítulo 2, secciones 2.4 y 2.5, proporciona una introducción rigurosa al formalismo de la traza parcial y su aplicación en sistemas cuánticos compuestos. El texto incluye numerosos ejemplos trabajados con estados de Bell y discute la conexión entre traza parcial y entrelazamiento. Especialmente relevante es la sección 2.4.3 sobre matrices densidad reducidas y su interpretación física.


\textbf{Preskill, J. \textit{Lecture Notes for Physics 229: Quantum Information and Computation}. California Institute of Technology, 2015}

Las notas del capítulo 2 sobre estados cuánticos y mediciones contienen una presentación muy clara de la traza parcial desde una perspectiva operacional. Se enfatiza la interpretación como «olvidar» información de un subsistema. Disponible gratuitamente en línea, incluye ejercicios resueltos sobre cálculo explícito de trazas parciales para estados multipartitos.


\textbf{Wilde, M. M. \textit{Quantum Information Theory}. Cambridge University Press, 2017}

El capítulo 3 presenta la teoría de operadores de densidad y traza parcial con un enfoque moderno orientado a la teoría de información cuántica. Particularmente útil es la sección 3.2 sobre purificaciones y la conexión entre estados mixtos y entrelazamiento con sistemas auxiliares. El libro desarrolla las propiedades de la traza parcial en el contexto de canales cuánticos, esencial para aplicaciones en corrección de errores y criptografía cuántica.


\textbf{Benenti, G., Casati, G., Strini, G. \textit{Principles of Quantum Computation and Information}. World Scientific, 2019}

El volumen I, capítulo 5, ofrece una introducción accesible a sistemas cuánticos compuestos con énfasis en ejemplos computacionales. Los autores incluyen código Mathematica para calcular trazas parciales de matrices densidad de dimensión arbitraria. La sección sobre mediciones parciales y decoherencia ilustra aplicaciones prácticas de la traza parcial en la descripción de sistemas abiertos.